\scanpage{50}
\uline{Ex 1} \unsure{$n$ fixé}, \unsure{prenons le foncteur} \unsure{donné par}
$\unsure{u_n^*(K)=K_n}$.
\[
M\longmapsto C^n(K_*,M)\overset{\unsure{\mathrm{d\acute{e}f}}}{=}
\operatorname{Hom}_{\unsure{\mathrm{Ens}}}(K_n,M)
=\operatorname{Hom}_{\unsure{\mathrm{Ens}}}(\unsure{u_n^*(K)},M)
\simeq \operatorname{Hom}_{\widehat{\Delta}}(K_*,\unsure{u_{n*}(M)}).
\]
\unsure{est représentable par} $\unsure{u_{n*}(M)}$.
\unsure{Puis, naturellement, les} $\unsure{u_{n*}(M)}$
forment un $k$-Module \unsure{cosimplicial} de $\widehat{\Delta}$, d'où
(\unsure{en passant au complexe} \illegible{3})
un complexe de \unsure{cochaînes} de $k$-Modules, qui est
$\unsure{C_{*M}^*}$. \unsure{D'où}
\[
C^*(K_*,M)\simeq\operatorname{Hom}(K_*,C_{*M}^*).
\]
(\unsure{Isom.}, \unsure{au choix}, de $k$-Modules \unsure{cosimpliciaux},
\unsure{ou de complexes de cochaînes}, \illegible{3} \unsure{le complexe} $K$).

On sait que pour $K_*=\Delta_n$, on a une \unsure{résolution}
de $M$, \unsure{donc} $C_{*M}^*$ est une résolution de $M_*$.
De plus les $\unsure{C_{*M}^n=u_{n*}(M)}$ sont \unsure{flasques},
\illegible{3}, $M$ \unsure{un} \illegible{1}!
\margin{NB $\unsure{C_{nM}^p=u_{p*}(M)_n
=\operatorname{Hom}(\Delta_n,\Delta_p)\cdot M
=M^{\operatorname{Hom}(\Delta_n,\Delta_p)}}$
$\unsure{=k^{\operatorname{Hom}(\Delta_n,\Delta_p)}\otimes_k M}$}

\uline{Ex 2} \unsure{Considérons le sous-complexe de cochaînes}
du $k$-Module \illegible{1} $C_{*M}^*$, formé des \unsure{cochaînes}
\unsure{non dégénérées}:
\[
\unsure{C_{*M}^{*!}=(C_{*M}^{n!})_{n\geq0}}:
\]
\[
\begin{aligned}
\unsure{C_{*M}^{n!}}
&=\bigcap_{\substack{\unsure{\sigma:\Delta_n\to\Delta_m}\\
\unsure{\sigma\text{ ép.},\,m<n}}}
\operatorname{Ker}(\sigma:C_{*M}^n\longrightarrow C_{*M}^m)\\
&=\bigcap_{\substack{\unsure{u_i:\Delta_n\to\Delta_{n-1}}\\
\unsure{i<n}}}
\operatorname{Ker}(\unsure{u_i}:C_{*M}^n\longrightarrow C_{*M}^{n-1})\\
&=\bigcap_{\substack{\unsure{\sigma:\Delta_n\to\Delta_m}\\
\unsure{m<n,\ \text{quelconque}}}}
\operatorname{Ker}(\sigma:C_{*M}^n\longrightarrow C_{*M}^m).
\end{aligned}
\]
On sait que cette \unsure{opération} \unsure{fournit} \illegible{4}
\unsure{un complexe} \illegible{2} \unsure{qui donne} \illegible{3}
\unsure{de} \illegible{3} \unsure{par} \illegible{2}
$\unsure{C^{*!}(X_*)\subset C^*(X_*)}$.
\margin{49}
